\chapter{博弈专题}\label{ux535aux5f08ux4e13ux9898}

本章只放竞赛中最常见的博弈模板。先判断游戏是否为无偏、有限、无环状态；能拆成独立子游戏时用 SG，双方都在完整状态上决策时用 Minimax。所有``必胜/必败''结论都要先确认终止规则和是否允许重复状态。

\section{Nim、Bash 与反常 Nim 取石子游戏}\label{nimbash-ux4e0eux53cdux5e38-nim-ux53d6ux77f3ux5b50ux6e38ux620f}

标准 Nim 每次从一堆中取任意正数个，所有堆异或和非零时先手必胜。Bash 游戏每次最多取 \passthrough{\lstinline!k!} 个，单堆 \passthrough{\lstinline!n!} 的必败位置是 \passthrough{\lstinline!n \% (k + 1) == 0!}。

\begin{lstlisting}[language={C++}]
// 接口：nim_first_win(piles)：返回标准 Nim 当前局面是否先手必胜；返回 true 表示本节条件成立/操作成功。
// 参数：各堆石子数 piles。
bool nim_first_win(const vector<int>& piles) {
    int x = 0; // 整个无偏 Nim 游戏的 SG 异或和。
    // 变量：v=当前相邻顶点/下一状态 v。
    for (int v : piles) x ^= v;
    return x != 0;
}

// 接口：bash_first_win(n, k)：返回 Bash 取石子局面是否先手必胜；返回 true 表示本节条件成立/操作成功。
// 参数：规模/上界 n；排名/选取数量 k。
bool bash_first_win(int n, int k) {
    // 一堆石子每次取 1..k；模 k+1 为 0 的位置是必败态。
    return n % (k + 1) != 0;
}
\end{lstlisting}

反常 Nim（最后取走者输）：若所有堆大小都为 \passthrough{\lstinline!1!}，先手在堆数为偶数时必胜；否则仍看异或和是否非零。

\begin{lstlisting}[language={C++}]
// 接口：misere_nim_first_win(piles)：返回反常 Nim 当前局面是否先手必胜；返回 true 表示本节条件成立/操作成功。
// 参数：各堆石子数 piles。
bool misere_nim_first_win(const vector<int>& piles) {
    // 只有“所有非空堆都是 1”时，最后一步输的规则会改变公式。
    // 变量：all_one=所有堆是否都只有 1 个石子的标记 all_one。
    bool all_one = true;
    // 变量：x=当前输入值/查询值/坐标 x；ones=石子数恰为 1 的堆数量 ones。
    int x = 0, ones = 0;
    // 变量：v=当前相邻顶点/下一状态 v。
    for (int v : piles) {
        all_one &= (v == 1);
        x ^= v;
        ones += (v == 1);
    }
    if (all_one) return ones % 2 == 0;
    return x != 0;
}
\end{lstlisting}

\section{通用 Sprague--Grundy 函数、最小未出现值与有向无环图博弈}\label{ux901aux7528-spraguegrundy-ux51fdux6570ux6700ux5c0fux672aux51faux73b0ux503cux4e0eux6709ux5411ux65e0ux73afux56feux535aux5f08}

无偏博弈的 Grundy 值满足 \passthrough{\lstinline!sg[u] = mex(\{sg[v]\})!}；独立子游戏的总值是异或。若状态图是 DAG，可以记忆化 DFS；若只是判断胜负，则``存在必败后继''为必胜，否则为必败。

\begin{lstlisting}[language={C++}]
// 接口：mex(values)：返回集合中没有出现的最小非负整数；返回类型 int。
// 参数：要求 mex 的非负整数集合 values。
int mex(vector<int> values) {
    // 排序去重后，从 0 开始找第一个断点；适合后继数量不太大的单次计算。
    sort(values.begin(), values.end());
    values.erase(unique(values.begin(), values.end()), values.end());
    // 变量：g=图/树的邻接表 g。
    int g = 0;
    // 变量：x=当前输入值/查询值/坐标 x。
    for (int x : values) {
        if (x == g) ++g;
        else if (x > g) break;
    }
    return g;
}

vector<int> sg;       // sg[u]=状态 u 的 Sprague-Grundy 值。
vector<char> sg_vis;  // sg_vis[u]=1 表示 sg[u] 已完成记忆化计算。
// 接口：grundy(u)：返回状态 x/u 的 Sprague-Grundy 值；返回类型 int。
// 参数：顶点/状态编号 u。
int grundy(int u) {
    // DAG 上记忆化求后继 SG；有环时必须换成专门的环博弈分析。
    if (sg_vis[u]) return sg[u];
    sg_vis[u] = true;
    // 变量：nxt=当前一步可到达的后继状态 nxt。
    vector<int> nxt;
    // 变量：v=当前边的终点 v。
    for (int v : dag[u]) nxt.push_back(grundy(v));
    return sg[u] = mex(nxt);
}

// 接口：first_win(starts)：返回若干独立子游戏异或和是否非零；返回 true 表示本节条件成立/操作成功。
// 参数：若干独立子游戏的起始状态 starts。
bool first_win(const vector<int>& starts) {
    // 只有子游戏互不影响时才能异或；共享资源的游戏不能套这里。
    // 变量：total=当前总数/连通块规模 total。
    int total = 0;
    // 变量：s=当前枚举的起点/源点 s。
    for (int s : starts) total ^= grundy(s);
    return total != 0;
}
\end{lstlisting}

\begin{lstlisting}[language={C++}]
vector<int8_t> win; // 0: 未求，1: 必胜，-1: 必败
// 接口：dag_game(u)：返回 DAG 状态 u 的 SG 值；返回类型 int。
// 参数：顶点/状态编号 u。
int dag_game(int u) {
    // 先暂定为必败，发现一个必败后继就改成必胜。
    if (win[u]) return win[u];
    win[u] = -1;
    // 变量：v=当前边的终点 v。
    for (int v : dag[u])
        if (dag_game(v) == -1) return win[u] = 1;
    return win[u];
}
\end{lstlisting}

\section{Nim 变体：取法集合与 Sprague--Grundy 周期}\label{nim-ux53d8ux4f53ux53d6ux6cd5ux96c6ux5408ux4e0e-spraguegrundy-ux5468ux671f}

如果单堆允许取的数量集合固定，先计算一维 SG；\passthrough{\lstinline!sg[n]!} 出现周期后可以直接按周期回答超大 \passthrough{\lstinline!n!}。多个互相独立的堆仍然异或。

\begin{lstlisting}[language={C++}]
// 接口：solve_subtraction(n, moves)：返回 n 堆/个石子的减法游戏 SG 值；返回类型 int。
// 参数：规模/上界 n；每步允许取走的石子数集合 moves。
int solve_subtraction(int n, const vector<int>& moves) {
    // 变量：sg=Sprague-Grundy 值数组 sg。
    vector<int> sg(n + 1);
    // 变量：seen=“Nim 变体：取法集合与 Sprague--Grundy 周期”这段代码中的临时量 seen；值由本行初始化式确定。
    vector<int> seen(moves.size() + 2);
    // 变量：x=当前输入值/查询值/坐标 x。
    for (int x = 1; x <= n; ++x) {
        // 变量：values=当前需要求 mex 的后继 SG 值集合 values。
        vector<int> values;
        // 变量：take=当前尝试取走的石子数量 take。
        for (int take : moves) if (take <= x)
            values.push_back(sg[x - take]);
        sg[x] = mex(values);
    }
    return sg[n];
}
\end{lstlisting}

\section{极大极小搜索（Minimax）与 Alpha-Beta 剪枝}\label{ux6781ux5927ux6781ux5c0fux641cux7d22minimaxux4e0e-alpha-beta-ux526aux679d}

双方零和、轮流行动且状态树有限时，Minimax 取当前方能保证的最优值。Alpha-Beta 剪枝维护祖先已知的下界 \passthrough{\lstinline!alpha!} 与上界 \passthrough{\lstinline!beta!}；排序优先搜索好着法会显著提高剪枝率。

\begin{lstlisting}[language={C++}]
// 接口：minimax(state, depth, alpha, beta, maximizing)：返回当前博弈状态在双方最优策略下的估值；返回类型 int。
// 参数：当前博弈状态 state；剩余搜索深度 depth；当前已知下界 alpha；当前已知上界 beta；当前层是否轮到极大化一方 maximizing。
int minimax(State& state, int depth, int alpha, int beta, bool maximizing) {
    // evaluate 始终站在固定一方视角；max 方取最大，min 方取最小。
    // alpha 是 max 已知下界，beta 是 min 已知上界。
    if (depth == 0 || terminal(state)) return evaluate(state);
    if (maximizing) {
        // 变量：best=目前找到的最优值 best。
        int best = INT_MIN;
        for (Move mv : ordered_moves(state)) {
            // apply/undo 必须成对，否则下一分支会继承上一分支的局面。
            apply(state, mv);
            best = max(best, minimax(state, depth - 1, alpha, beta, false));
            undo(state, mv);
            alpha = max(alpha, best);
            // beta cutoff：min 方已有更小上界，后面的分支不会被选择。
            if (alpha >= beta) break; // beta cutoff
        }
        return best;
    }
    // 变量：best=目前找到的最优值 best。
    int best = INT_MAX;
    for (Move mv : ordered_moves(state)) {
        apply(state, mv);
        best = min(best, minimax(state, depth - 1, alpha, beta, true));
        undo(state, mv);
        beta = min(beta, best);
        // alpha cutoff：max 方已有更大下界，后面的分支不会被选择。
        if (alpha >= beta) break; // alpha cutoff
    }
    return best;
}
\end{lstlisting}

实际写棋类题时，\passthrough{\lstinline!State/Move/apply/undo/terminal/evaluate!} 替换为题目状态；有深度限制时应加入深度奖励、置换表和走法排序，避免把``未搜到底''的局面误判成终局。

\section{博弈题检查清单}\label{ux535aux5f08ux9898ux68c0ux67e5ux6e05ux5355}

\begin{itemize}
\tightlist
\item
  终止状态是``不能走''输，还是``最后一步''输？这决定普通/反常规则。
\item
  是否无偏？若双方可用的操作不同，不能直接套 SG xor。
\item
  状态是否可拆成独立和？有共享资源时不能强行异或。
\item
  状态图是否有环？有环时需要周期、记忆化博弈或胜负状态方程。
\item
  题目只问胜负用 Win/Lose，问组合价值或多堆合并才引入 Grundy。
\end{itemize}
